\section*{TÓM TẮT}
\phantomsection\addcontentsline{toc}{section}{\numberline {} TÓM TẮT}
Đồ án chuyên ngành I tập trung nghiên cứu lý thuyết và kiến trúc của các hệ thống tính toán song song phân tán, đóng vai trò cốt lõi trong việc huấn luyện và suy luận các mô hình Trí tuệ Nhân tạo (AI) quy mô lớn hiện nay. Nội dung báo cáo được tổ chức thành bốn phần chính, đi từ các khái niệm cơ bản về lập trình song song trên CPU với OpenMP và MPI, đến việc khai thác sức mạnh của kiến trúc phần cứng GPU thông qua lập trình CUDA. Tiếp đó, báo cáo phân tích sâu các cơ chế giao tiếp liên GPU tiên tiến như NCCL và NVSHMEM, vốn là nền tảng cho việc luân chuyển dữ liệu ở tốc độ cao. Trọng tâm của đồ án nằm ở việc khảo sát các chiến lược song song hóa (Data Parallelism, Tensor Parallelism, Pipeline Parallelism, Context Parallelism) và cách chúng được hiện thực hóa trong các framework hàng đầu như PyTorch Distributed, DeepSpeed, và Megatron Core. Cuối cùng, sự khác biệt về yêu cầu phần cứng, bộ nhớ và băng thông mạng giữa hai giai đoạn huấn luyện (training) và suy luận (inference) cũng được làm rõ. Đồ án cung cấp một cái nhìn toàn cảnh, tạo cơ sở vững chắc cho việc triển khai và tối ưu hóa các hệ thống AI phân tán trong thực tế.

\newpage
\cleardoublepage